\section*{第 \ref{ch:g12:stereochemistry} 章——立体化学：手性与异构体}

\begin{solution}{exo:g12:stereochemistry:1}
2-丁醇：碳 2（\ce{H}、\ce{OH}、\ce{CH3}、\ce{CH2CH3}）。
2-丙醇：无（碳 2 连有两个相同的 \ce{CH3}）。
2-氯丁烷：碳 2。乳酸：碳 2（\ce{H}、\ce{OH}、
\ce{CH3}、\ce{COOH}）。
\end{solution}

\begin{solution}{exo:g12:stereochemistry:2}
第一个：两个 \ce{Cl}（优先级高于 \ce{H}）位于同侧：Z。
第二个：左侧的 \ce{CH3} 优先级更高，位于上侧；右侧的
\ce{CH2CH3} 优先级更高，位于下侧：E。
\end{solution}

\begin{solution}{exo:g12:stereochemistry:3}
有手性：手套、螺丝钉。无手性：勺子、立方体（每个都有
一个镜面）。\ce{CHFClBr} 有手性（碳原子上连有四个不同的原子）；
\ce{CH2Cl2} 无手性。
\end{solution}

\begin{solution}{exo:g12:stereochemistry:4}
等量两种对映异构体的混合物。外消旋香芹酮含有
两种对映异构体，因此鼻子同时接收两种气味。
\end{solution}

\begin{solution}{exo:g12:stereochemistry:5}
否：碳 1 连有两个相同的原子，即 \ce{H} 和 \ce{H}；交换
它们所得仍为同一分子。
\end{solution}

\begin{solution}{exo:g12:stereochemistry:6}
\setchemfig{atom sep=2em}
\chemfig{C(-[2]CH_2CH_3)(-[4]H_3C)(<[:-30]OH)(<:[:-150]H)} \quad 镜像
\quad \chemfig{C(-[2]CH_2CH_3)(-[0]CH_3)(<[:-150]HO)(<:[:-30]H)}
\end{solution}

\begin{solution}{exo:g12:stereochemistry:7}
2-氯丁烷：2 种（一个手性碳原子）。2,3-二氯丁烷：3 种（两个
手性碳原子，两半相同：一对对映异构体与一个内消旋
体）。2-戊烯：2 种（Z 和 E）。
\end{solution}

\begin{solution}{exo:g12:stereochemistry:8}
对映异构体；非对映异构体；非对映异构体；同一分子（内消旋
酒石酸与其镜像可以重合）。
\end{solution}

\begin{solution}{exo:g12:stereochemistry:9}
在交错构象中，两个 \ce{CH3} 基团位于相反方向，大基团
彼此尽可能远离：该构象最稳定。在
重叠构象中，两个 \ce{CH3} 相对，彼此拥挤。
\end{solution}

\begin{solution}{exo:g12:stereochemistry:10}
其两半各为 \ce{CH(OH)COOH}，二者互为镜像：一个
平面位于两个中心碳原子之间，将分子切成两个互为镜像的
半部分。具有镜面的分子与其镜像相同：
无手性。
\end{solution}

\begin{solution}{exo:g12:stereochemistry:11}
\setchemfig{atom sep=1.7em}
(Z)-2-戊烯 \chemfig{C(-[:150]H_3C)(-[:210]H)=C(-[:30]CH_2-CH_3)(-[:-30]H)}
和 (E)-2-戊烯 \chemfig{C(-[:150]H_3C)(-[:210]H)=C(-[:30]H)(-[:-30]CH_2-CH_3)}。
\end{solution}

\begin{solution}{exo:g12:stereochemistry:12}
与酒石酸类似，只是用 \ce{Br} 代替 \ce{OH}，\ce{CH3}
代替 \ce{COOH}：两个 \ce{Br} 均为实楔形，均为虚楔形（一对
对映异构体），以及一实一虚，后者有镜面：即
内消旋体。共三种立体异构体。
\end{solution}

\begin{solution}{exo:g12:stereochemistry:13}
一半。化学家可以拆分两种对映异构体（如巴斯德所做的，用的是
其他方法），或者只合成有用的对映异构体，采用
本身具有手性的反应（手性催化剂、酶或手性起始
原料）。
\end{solution}

\begin{solution}{exo:g12:stereochemistry:14}
三种：(2E,4E)、(2Z,4Z) 和 (2E,4Z)，其中第三种与从另一端读取的
(2Z,4E) 为同一分子。
\end{solution}

\begin{solution}{exo:g12:stereochemistry:15}
$2^3 = 8$ 种立体异构体，即 4 对对映异构体。每种
立体异构体有一个对映异构体，并有 $8 - 2 = 6$ 个非对映异构体。
\end{solution}

\begin{solution}{pb:g12:stereochemistry:1}
\textbf{1.} \ce{C4H6O6}；$M = 4 \times 12.0 + 6 \times 1.0 + 6 \times 16.0
= \qty{150.0}{g/mol}$。

\textbf{2.} 两个羧基（酸）和两个羟基
（醇）。

\textbf{3.} 两个中心碳原子，各与 \ce{H}、\ce{OH}、
\ce{COOH} 和 \ce{CH(OH)COOH} 相连。

\textbf{4.} $2^2 = 4$。

\textbf{5.} 两者均为实楔形。

\textbf{6.} 两个 \ce{OH} 均为虚楔形；它不能与
L 重合：它是 D-酒石酸，即其对映异构体。

\textbf{7.} 绕中心键旋转，使两个 \ce{OH} 朝向
同一方向，分子在两个中心碳原子之间有一个平面，
该平面将一半反射到另一半上。

\textbf{8.} 否：它是内消旋酒石酸。

\textbf{9.} 在四种组合（实，实）、（虚，虚）、
（实，虚）、（虚，实）中，其中后两种为同一分子，
即内消旋形式，因为分子的两半相同。

\textbf{10.} L 与 D：对映异构体。L 与内消旋体：非对映异构体。

\textbf{11.} 在 \qty{169}{\celsius} 时同样如此：对映异构体具有相同的
物理性质。

\textbf{12.} 否：它是 L 的非对映异构体，是
另一种具有自身性质的化合物。

\textbf{13.} 鼻子中的受体具有手性，能够区分
互为镜像的分子；温度计没有手性，做不到这一点。

\textbf{14.} 否，这是两种化合物的混合物；整体上不具有
光学活性，两种对映异构体的旋光作用相互抵消。

\textbf{15.} 只含一种对映异构体：每颗晶体都
仅由 L 或仅由 D 构成。

\textbf{16.} 各 \qty{5.0}{g}。

\textbf{17.} 内消旋体是另一种化合物，不属于
L 与 D 的外消旋混合物。

\textbf{18.} 将每堆晶体分别溶解：两种溶液使偏振光
发生相同角度、方向相反的旋转（两种晶体
互为镜像）。

\textbf{19.} 共四种：其两个手性碳原子所连基团不同
（一个连 \ce{Cl}，另一个连 \ce{OH}），因此没有任何组合具有对称面，
且任意两种组合都不是同一分子。

\textbf{20.} 三种：L、D 与内消旋体。
\end{solution}
